% Section 19.3: Statement and proof | Part: 3 | Depends on: sec:stokes, sec:interior, sec:metric | Status: done | Original lines: 2865-2891
%!TEX root = ../../../main.tex
\subsection{Statement and proof}

\begin{theorem}[Gauss--Ostrogradsky, divergence theorem]\label{thm:gauss}
Let $(M,g)$ be a compact oriented Riemannian $n$-manifold with boundary and $X$ a smooth vector field. Then
\[
\int_M\operatorname{div}X\,\dd V=\int_{\partial M}\langle X,N\rangle\,\dd A .
\]
\end{theorem}
\begin{proof}
By \eqref{eq:divmu}, $\operatorname{div}X\,\vol_g=\dd(\iota_X\vol_g)$. Apply Theorem~\ref{thm:stokes} to $\omega=\iota_X\vol_g$ and use Lemma~\ref{lem:fluxbdry}.
\end{proof}

\begin{example}[The position field on the unit ball]\label{ex:ball}
Let $X=x\partial_x+y\partial_y+z\partial_z$ on the unit ball $B\subset\R^3$.
\begin{enumerate}
\item \emph{Volume side.} $\operatorname{div}X=\partial_xx+\partial_yy+\partial_zz=3$, so $\int_B3\,\dd V=3\cdot\frac{4\pi}3=4\pi$.
\item \emph{Boundary side, with the formula of the lemma.} On the unit sphere the outward normal is $N=(x,y,z)=X$, hence $\langle X,N\rangle=x^2+y^2+z^2=1$, and the flux is $\int_{S^2}1\,\dd A=4\pi$.
\item \emph{Boundary side, directly with forms.} $\iota_X\vol=x\,\dd y\wedge\dd z+y\,\dd z\wedge\dd x+z\,\dd x\wedge\dd y$ (formula \eqref{eq:iotavol}). At the north pole $(0,0,1)$ the tangent plane is spanned by $\partial_x,\partial_y$, where $\dd z=0$, so the restriction is $z\,\dd x\wedge\dd y=\dd x\wedge\dd y=\dd A$, again with coefficient $1$.
\end{enumerate}
Both sides equal $4\pi$ \checkmark.
\end{example}

\begin{synthesis}[the area form of the sphere is a contraction]\label{syn:sphere}
In Section~\ref{sec:tensorfields} the area form of the sphere of radius $r$ was defined by $\alpha_p(v_1,v_2)=\frac1r(v_1\times v_2)\cdot p$. Since $(v_1\times v_2)\cdot N=\vol(N,v_1,v_2)$ for $N=p/r$ (triple product),
$\alpha=\iota_N\vol$ restricted to the sphere: this is Lemma~\ref{lem:fluxbdry} with $X=N$, $\langle N,N\rangle=1$. In spherical coordinates it is $r^2\sin\theta\,\dd\theta\wedge\dd\varphi$, with total area $4\pi r^2$, and Example~\ref{ex:ball} (flux of the position field through $S^2$, $=4\pi$) is the case $r=1$.
\end{synthesis}
